
\documentclass[11pt,a4paper]{article}
\usepackage[margin=1in]{geometry}
\usepackage{amsmath,amssymb,amsthm,mathtools}

\title{Merged Solutions: Exercises 1.1--1.3, 1.5--1.7 and 2.1--2.8, 3.1}
\author{}
\date{\today}

\begin{document}
\maketitle

\section*{Exercise 1.1}
\begin{enumerate}
\renewcommand\labelenumi{(\alph{enumi})}
\item Show that the inner product has the following properties for all $x,y,z\in\mathbb{R}^n$ and $a\in\mathbb{R}$:
\[
\langle x,y\rangle=\langle y,x\rangle,\qquad
\langle x+y,z\rangle=\langle x,z\rangle+\langle y,z\rangle,\qquad
\langle ax,y\rangle=a\langle x,y\rangle.
\]
\textit{Solution.} These follow from the Euclidean inner product $\langle x,y\rangle=\sum_{i=1}^n x_i y_i$.

\item For $t\in\mathbb{R}$ and $x,y\in\mathbb{R}^n$ show
\[
\|x+ty\|^2=\|x\|^2+2t\langle x,y\rangle+t^2\|y\|^2\ge 0.
\]
\textit{Solution.} Expand $\langle x+ty,x+ty\rangle$.

\item Deduce Cauchy--Schwarz: $|\langle x,y\rangle|\le \|x\|\,\|y\|$.\\
\textit{Solution.} View $\|x+ty\|^2$ as a quadratic in $t$ and use non-positivity of its discriminant.
Equality iff $x,y$ are linearly dependent.

\item Deduce the triangle inequality $\|x+y\|\le \|x\|+\|y\|$.

\item Show the reverse triangle inequality $|\|x\|-\|y\||\le \|x-y\|$.

\item Suppose $x=(x^1,\dots,x^n)^t$.
\begin{enumerate}
\renewcommand\labelenumii{(\roman{enumii})}
\item $\displaystyle \max_{k}|x^k|\le \|x\|$.
\item $\displaystyle \|x\|\le \sqrt{n}\,\max_{k}|x^k|$.
\end{enumerate}
\textit{Solution.} Immediate from $\|x\|^2=\sum (x^i)^2$.
\end{enumerate}

\section*{Exercise 1.2}
Let $(x_i)$ and $(y_i)$ in $\mathbb{R}^n$ satisfy $x_i\to x$, $y_i\to y$.
\begin{enumerate}
\renewcommand\labelenumi{(\alph{enumi})}
\item $x_i+y_i\to x+y$. \textit{Proof.} $\|(x_i+y_i)-(x+y)\|\le \|x_i-x\|+\|y_i-y\|\to 0$.
\item $\langle x_i,y_i\rangle\to \langle x,y\rangle$; in particular $\|x_i\|\to \|x\|$.\\
\textit{Proof.} Expand $\langle x_i,y_i\rangle-\langle x,y\rangle$ and use Cauchy--Schwarz.
\item If $a_i\to a$ in $\mathbb{R}$, then $a_i x_i\to ax$.\\
\textit{Proof.} $a_ix_i-ax=(a_i-a)(x_i-x)+(a_i-a)x+a(x_i-x)$.
\end{enumerate}

\section*{Exercise 1.3}
Which of the following subsets of $\mathbb{R}^n$ is open?
\begin{enumerate}
\renewcommand\labelenumi{(\alph{enumi})}
\item $\mathbb{R}^n$: \textbf{Yes}.
\item $\varnothing$: \textbf{Yes}.
\item $\{x\in\mathbb{R}^n:x^1>0\}$: \textbf{Yes}.
\item $\{x\in\mathbb{R}^n:x^i\in[0,1]\}$: \textbf{No}.
\item $\mathbb{Q}^n$: \textbf{No}.
\end{enumerate}

\section*{Exercise 1.5}
\begin{enumerate}
\renewcommand\labelenumi{(\alph{enumi})}
\item If $x_i\to x$ and $\|x_i\|<r$ for all $i$, then $\|x\|\le r$.\\
\textit{Proof.} Contradict $\|x\|>r$ using $\varepsilon=(\|x\|-r)/2$ and the reverse triangle inequality.
\item The closure of $B_r(y)=\{x:\|x-y\|<r\}$ is $\overline{B_r(y)}=\{x:\|x-y\|\le r\}$.\\
\textit{Proof.} Use (a) for the first inclusion; for boundary points take $x_i=y+(1-\tfrac1i)(x-y)$.
\end{enumerate}

\section*{Exercise 1.6}
Find $A^\circ$, $\overline{A}$ and $\partial A$.
\begin{enumerate}
\renewcommand\labelenumi{(\alph{enumi})}
\item $A=\mathbb{R}^n$: $A^\circ=\mathbb{R}^n$, $\overline{A}=\mathbb{R}^n$, $\partial A=\varnothing$.
\item $A=\{0\}$: $A^\circ=\varnothing$, $\overline{A}=\{0\}$, $\partial A=\{0\}$.
\item $A=\{x:x^1\ge 0\}$: $A^\circ=\{x:x^1>0\}$, $\overline{A}=A$, $\partial A=\{x:x^1=0\}$.
\item $A=\{x:x^i\in[0,1]\ \forall i\}$: $A^\circ=\{x:x^i\in(0,1)\ \forall i\}$, $\overline{A}=A$, $\partial A=A\setminus A^\circ$.
\item $A=\{x:x^i\in\mathbb{Q}\ \forall i\}$: $A^\circ=\varnothing$, $\overline{A}=\mathbb{R}^n$, $\partial A=\mathbb{R}^n$.
\end{enumerate}

\section*{Exercise 1.7}
\begin{enumerate}
\renewcommand\labelenumi{(\alph{enumi})}
\item If $U_1,U_2$ are open, then $U_1\cup U_2$ and $U_1\cap U_2$ are open.
\item If $E_1,E_2$ are closed, then $E_1\cup E_2$ and $E_1\cap E_2$ are closed.
\item For any collection $\mathcal U$ of open sets: $\bigcup\mathcal U$ is open; $\bigcap\mathcal U$ need not be open (e.g.\ $\cap_{n\ge1}(-1/n,1/n)=\{0\}$).
For closed sets: arbitrary intersections are closed, finite unions are closed, infinite unions need not be closed.
\end{enumerate}

\section*{Exercise 2.1}
A sequence $(x_i)\subset\mathbb{R}^n$ is Cauchy if $\forall\varepsilon>0\,\exists N:\ i,j\ge N\Rightarrow \|x_i-x_j\|<\varepsilon$.
\begin{enumerate}
\renewcommand\labelenumi{(\alph{enumi})}
\item If $x_i\to x$, then $(x_i)$ is Cauchy. \textit{Proof.} Triangle inequality with $\varepsilon/2$.
\item If $(x_i)$ is Cauchy, then it is bounded. \textit{Proof.} Take $\varepsilon=1$.
\item Every Cauchy sequence in $\mathbb{R}^n$ converges. \textit{Proof.} Bounded $\Rightarrow$ has a convergent subsequence; use Cauchy property to pass to the full sequence.
\end{enumerate}
\textbf{Remark.} In $\mathbb{R}$, $\mathbb{R}^n$, and $\mathbb{C}\cong\mathbb{R}^2$, a sequence is convergent iff it is Cauchy (completeness). In incomplete spaces (e.g.\ $\mathbb{Q}$), Cauchy does not imply convergent.

\section*{Exercise 2.2}
Suppose $A\subset\mathbb{R}^n$ and $f:A\to\mathbb{R}^m$. Then
\[
\lim_{x\to p}f(x)=F \iff \text{ for every sequence } x_i\in A,\ x_i\neq p,\ x_i\to p,\ \text{we have } f(x_i)\to F.
\]
\textit{Proof.} The forward direction is immediate from the $\varepsilon$--$\delta$ definition; the reverse is proved by contradiction constructing a sequence $x_k\to p$ with $f(x_k)$ staying $\varepsilon_0$ away from $F$.

\section*{Exercise 2.3}
\begin{enumerate}
\renewcommand\labelenumi{(\alph{enumi})}
\item $f:\mathbb{R}\to\mathbb{R}^n$, $x\mapsto (x,0,\dots,0)^t$ is continuous: $\|f(x)-f(x_0)\|=|x-x_0|$.
\item For $A\subset\mathbb{R}^n$, $f=(f^1,\dots,f^m):A\to\mathbb{R}^m$ is continuous at $p$ iff each $f^k$ is continuous at $p$.
\item Any polynomial map $f:\mathbb{R}^n\to\mathbb{R}$ is continuous (built from projections via sums/products).
\end{enumerate}

\section*{Exercise 2.4}
Let $E\subset\mathbb{R}^n$ be compact and $f:E\to\mathbb{R}^m$ continuous.
\begin{enumerate}
\renewcommand\labelenumi{(\alph{enumi})}
\item $f(E)$ is bounded. \textit{Proof.} Each component $f^j$ attains its max/min on $E$.
\item $f(E)$ is closed, hence compact. \textit{Proof.} Use sequential closedness via compactness of $E$.
\item If $E$ is merely closed, $f(E)$ need not be closed; e.g.\ $f(x)=\arctan x$ on $E=\mathbb{R}$.
\end{enumerate}

\section*{Exercise 2.5 (*)}
\begin{enumerate}
\renewcommand\labelenumi{(\alph{enumi})}
\item If $f:\mathbb{R}^n\to\mathbb{R}^m$ is continuous and $U$ open in $\mathbb{R}^m$, then $f^{-1}(U)$ is open.
\item If $f^{-1}(U)$ is open for every open $U\subset\mathbb{R}^m$, then $f$ is continuous on $\mathbb{R}^n$.
\end{enumerate}

\section*{Exercise 2.6}
Assume $f:\mathbb{R}\to\mathbb{R}$ differentiable with $|f'(x)|\le M$ for all $x$.
\begin{enumerate}
\renewcommand\labelenumi{(\alph{enumi})}
\item $|f(x)-f(y)|\le M|x-y|$ by MVT.
\item Hence $f$ is uniformly continuous (Lipschitz with constant $M$).
\end{enumerate}

\section*{Exercise 2.7}
Which functions are (i) continuous, (ii) uniformly continuous on their domains?
\begin{enumerate}
\renewcommand\labelenumi{(\alph{enumi})}
\item $f:\mathbb{R}\to\mathbb{R}$, $x\mapsto e^x$: continuous, not uniformly continuous.
\item $f:(0,1)\to\mathbb{R}$, $x\mapsto e^x$: continuous, uniformly continuous (MVT and $e^c\le e$).
\item $f:\mathbb{R}\to\mathbb{R}$, $x\mapsto\sin x$: continuous, uniformly continuous ($|\sin x-\sin y|\le |x-y|$).
\item $f:[0,\infty)\to\mathbb{R}$, $x\mapsto \sqrt{x}$: continuous, uniformly continuous ($|\sqrt{x}-\sqrt{y}|\le \sqrt{|x-y|}$).
\item[\textit{f)}] $f:\mathbb{R}^n\setminus\{0\}\to\mathbb{R}^n$, $x\mapsto x/\|x\|$: continuous, not uniformly continuous (fails near $0$).
\end{enumerate}

\section*{Exercise 2.8}
Let $A\subset \mathbb{R}^n$ and $f,g:A\to\mathbb{R}$.
\begin{enumerate}
\renewcommand\labelenumi{(\alph{enumi})}
\item Suppose $f$ and $g$ are uniformly continuous and bounded on $A$. Then $fg$ is uniformly continuous and bounded on $A$.

\textit{Solution (boundedness).} If $|f|\le M_f$ and $|g|\le M_g$ on $A$, then $|fg|\le M_fM_g$.

\textit{Solution (uniform continuity).} Fix $\varepsilon>0$ and set
$\eta=\min\!\left(1,\frac{\varepsilon}{2(M_f+M_g+1)}\right)$.
By uniform continuity of $f$ and $g$, there exists $\delta>0$ such that
$\|x-y\|<\delta$ implies $|f(x)-f(y)|<\eta$ and $|g(x)-g(y)|<\eta$.
Then for such $x,y$,
\begin{align*}
|f(x)g(x)-f(y)g(y)|
&= |[f(x)-f(y)][g(x)-g(y)] + g(y)[f(x)-f(y)] + f(y)[g(x)-g(y)]|\\
&\le \eta^2 + M_g\,\eta + M_f\,\eta
\;\le\; (M_f+M_g+1)\eta \;\le\; \varepsilon.
\end{align*}
Hence $fg$ is uniformly continuous.

\item If $f$ and $g$ are uniformly continuous but not necessarily bounded, $fg$ need not be uniformly continuous.

\textit{Counterexample.} On $\mathbb{R}$, take $f(x)=x$ and $g(x)=\sin x$. Both are uniformly continuous (Lipschitz).
Assume $x\sin x$ were uniformly continuous. For $\varepsilon=1$, let $\delta>0$ be its modulus.
Set $t=\min(\delta/2,1)$ and take $y_k=k\pi$, $x_k=k\pi+t$.
Then $|x_k-y_k|=t<\delta$ but
\[
|x_k\sin x_k - y_k\sin y_k| = |(k\pi+t)\sin t| \ge k\pi \sin t \xrightarrow{k\to\infty} \infty,
\]
contradicting uniform continuity. Thus $x\sin x$ is not uniformly continuous.
\end{enumerate}

\section*{Exercise 3.1}
For $i\in\mathbb{N}$, define $f_i:[0,1]\to\mathbb{R}$ by
\[
f_i(x)=
\begin{cases}
2^i x, & 0\le x<2^{-i},\\[2pt]
2-2^i x, & 2^{-i}\le x<2^{-(i-1)},\\[2pt]
0, & x\ge 2^{-(i-1)}.
\end{cases}
\]

\begin{enumerate}
\renewcommand\labelenumi{(\alph{enumi})}
\item \textit{Sketch.} $f_i$ is a triangular ``spike'' supported on $[0,2^{-(i-1)}]$: it rises linearly
from $0$ at $x=0$ to height $1$ at $x=2^{-i}$, then decreases linearly back to $0$ at $x=2^{-(i-1)}$; afterwards it is identically $0$.
In particular, $\displaystyle \sup_{x\in[0,1]} f_i(x)=1$ for every $i$.

\item \textit{Pointwise limit.} For any fixed $x>0$ choose $I$ such that $2^{-(i-1)}<x$ for all $i\ge I$; then $f_i(x)=0$ for $i\ge I$.
Also $f_i(0)=0$ for all $i$. Hence $f_i(x)\to 0$ for every $x\in[0,1]$.

\item \textit{Interchange of $\lim$ and $\sup$.} We have
\[
\lim_{i\to\infty}\ \sup_{x\in[0,1]} f_i(x}
= \lim_{i\to\infty} 1 = 1,
\qquad\text{while}\qquad
\sup_{x\in[0,1]} \ \lim_{i\to\infty} f_i(x) = \sup_{x\in[0,1]} 0 = 0.
\]
Therefore the equality
\[
\lim_{i\to\infty} \sup_{x\in[0,1]} f_i(x) = \sup_{x\in[0,1]} \lim_{i\to\infty} f_i(x)
\]
is \emph{false} for this family. (This also shows $\{f_i\}$ does not converge uniformly to $0$.)
\end{enumerate}

\end{document}
